%!TEX root = ../LLMSeminar.tex
% ──────────────────────────────────────────────────────────────
%  Section 3 — From Function to Agent
% ──────────────────────────────────────────────────────────────

\section{From function to agent}\label{sec:agents}

\subsection{Two ways in}

There are two ways to interact with an LLM\@.  The first is through a
web interface---ChatGPT, Claude.ai---where the chat UI hides the
underlying API\@.  The second is through a terminal: cURL, Python, or a
coding agent like Claude Code.  Both use the same API underneath.

The web interface is convenient for casual queries, but it hides the
mechanism.  To \emph{use} an LLM effectively for research, you need to
understand and control that mechanism.

\subsection{Coding agents}

The critical development of 2025 was the emergence of \textbf{coding
agents}: programs that empower LLMs with autonomous tool use.  Claude
Code, Cursor, and Cline are the leading examples.

The idea is simple.  Instead of a human reading the LLM's output and
manually executing suggestions, a \emph{scaffolding program} does it
automatically in a loop.

\subsection{The agent loop}

\begin{keyresult}[The agent loop]
  The complete architecture of every coding agent is:
  \begin{lstlisting}
while not done:
    response = llm(system_prompt + history)
    tool_calls = parse(response)
    results = execute(tool_calls)
    history += [response, results]
  \end{lstlisting}
  That is it.  Claude Code, Cursor, Copilot---they all implement this
  loop with more tools and better prompts, but the architecture is
  identical.
\end{keyresult}

Let us unpack each step:

\begin{enumerate}[(1)]
  \item The scaffolding sends the system prompt and accumulated history
    to the LLM as a single function call.
  \item The LLM responds with text that may include structured
    \emph{tool-call requests}: ``please read file \code{X}'' or ``please
    execute command~\code{Y}.''
  \item The scaffolding parses these requests and executes them locally.
  \item The results are appended to the history, and the loop continues.
\end{enumerate}

\begin{center}
\begin{tikzpicture}[node distance=1.4cm, every node/.style={align=center}]
  \node[sysbox, fill=spacecadet, text=white, minimum width=4cm,
        minimum height=0.9cm, font=\small]
    (sys) {System prompt + History};

  \node[sysbox fill, below=1cm of sys, minimum width=4cm,
        minimum height=0.9cm, font=\small]
    (llm) {\textbf{LLM}\; {\scriptsize $f(s) \to s$}};

  \node[sysbox, right=2.5cm of llm, minimum height=0.9cm, font=\small]
    (parse) {Parse tool calls};

  \node[sysbox, below=1cm of parse, minimum height=0.9cm, font=\small]
    (exec) {Execute tools};

  \draw[sysarrow] (sys) -- (llm);
  \draw[sysarrow] (llm) -- (parse)
    node[midway, above, font=\scriptsize] {response};
  \draw[sysarrow] (parse) -- (exec)
    node[midway, right, font=\scriptsize] {call};
  \draw[sysarrow] (exec.south) -- ++(0,-0.5) -| ([xshift=-0.4cm]llm.south)
    node[near start, below, font=\scriptsize] {results};
  \draw[sysarrow, dashed, color=munsell]
    ([xshift=0.4cm]llm.north) -- ([xshift=0.4cm]sys.south)
    node[midway, right, font=\scriptsize, text=munsell] {append};
\end{tikzpicture}
\end{center}

\subsection{Tools}

You describe tools in the system prompt.  The LLM generates structured
text requesting a tool call.  Your scaffolding parses and executes it.

\begin{keyresult}[The key insight]
  The LLM never executes anything directly.  Your scaffolding acts on
  its text.  The LLM is an oracle that emits instructions; the
  scaffolding is the hands.
\end{keyresult}

Even two primitive tools suffice to do useful work:

\begin{lstlisting}
read_file(path)           -> contents
write_file(path, content) -> ok
\end{lstlisting}

With just these two tools, an agent can read existing code, understand
it, write new code, and iterate.  Production agents like Claude Code
have many more tools---bash execution, web search, file
manipulation---but the principle is the same.

\begin{greybox}[Further reading]
  Thorsten Ball, \emph{How to Build an Agent}---ampcode.com.  A clear,
  practical guide to the agent loop.
\end{greybox}

\subsection{The feedback loop}

\begin{audience}[Audience question]
  \emph{Is writing a good initial prompt the most important thing?}

  No.  Prompt engineering is dead.  The key to improving an LLM's output
  is \textbf{feedback}.  Tighten the feedback loop.  Tell it what is
  wrong so it can correct itself and try again.  The magic is not in the
  prompt; it is in the iteration.

  How does it correct itself?  It re-inputs the entire history---including
  your correction---back into the function.
\end{audience}

\subsection{Live demonstration}

At this point in the talk, we launched Claude Code live and asked an
audience member for a research problem.  The challenge: compute the
potential energy curves for approximately 25 electronic states of FeO
(iron oxide).  Claude Code proceeded to:

\begin{enumerate}[(1)]
  \item Survey the system for available quantum-chemistry software.
  \item Find none installed.
  \item Propose to install PySCF (a Python-based quantum-chemistry
    package).
  \item Set up a virtual environment.
  \item Write a smoke test for a small active space.
  \item Design a production computation script using SA-CASSCF(12,12) +
    SC-NEVPT2, with X2C relativistic corrections, a cc-pVTZ-DK basis
    set, 35 bond lengths, and ${\sim}25$ electronic states.
  \item Execute the computation in the background while the talk
    continued.
\end{enumerate}

The entire process was autonomous.  The only human input was the initial
problem statement and a ``yes'' when Claude asked permission to install
software.
